\documentclass[12pt]{article}
\usepackage[a4paper,margin=1in]{geometry}
\usepackage{amsmath,amssymb,mathtools}
\usepackage{fontspec}
\usepackage{polyglossia}
\setmainlanguage{hebrew}
\setotherlanguage{english}
\newfontfamily\hebrewfont{Arial}
\setlength{\parindent}{0pt}
\setlength{\parskip}{0.6em}

\begin{document}
% OCR draft generated by MathGrade.
% Source: אופק סגל_13030568_assignsubmission_file_תרגיל 1 חדוא.pdf
% Review carefully before grading.
% Required grading markers:
%   \subsection*{Question 1}
%   \textbf{(א)}
%   \textbf{(ב)}

% WARNING: MathGrade could not confidently detect question titles.
% Before grading, split this OCR text into:
% \subsection*{Question 1}
% \textbf{(א)}
% \textbf{(ב)}

\subsection*{Question 1}

\[
|a| \leq 5
\]
\[
\begin{aligned}
& |x-2| \leq 5 \quad k \\
-5 \leq x-2 \leq 5 \quad & 1+2 \\
-3 \leq x \leq 7 &
\end{aligned}
\]
\[
\begin{aligned}
& |x+20| \leq 2 \\
& -2 \leq x+20 \leq 2 \\
& -22 \leq x \leq-18
\end{aligned}
\]
\[
\begin{gathered}
-1>y-6 \\
s>y
\end{gathered}
\]
\[
\begin{array}{ll}
\text { II } & 1<y-6 \\
k & 7<y
\end{array}
\]

\[
\begin{gathered}
x<-3: 3 \\
-(x+2)-(x+3) \leq 2 \\
-x-2-x-3 \leq 2 \\
-2 x-5 \leq 2 \\
-2 x \leq 7 \\
x \geq-3.5 \\
-3.5<x<-3
\end{gathered}
\]
\[
\begin{gathered}
-3 \leq x \leq-2 \\
-(x+2)+x+3 \leq 2 \\
-x-2+x+3 \leq 2 \\
1 \leq 2 \\
-3 \leq x \leq-2
\end{gathered}
\]
\[
\begin{gathered}
-2<x: 1 \\
x+2 \rightarrow x+3 \leq 2 \\
2 x \leq-3 \\
x \leq-1.5 \\
-2<x \leq-1.5
\end{gathered}
\]

معاد.
\[
-3.5 \leq x \leq-1.5
\]
\[
|x+5|+|x+7|>1
\]
\[
x=-7 \quad-1 \quad x=-5 \quad-2 \text { ork }
\]
\[
\begin{aligned}
& \text { : } 3 \text { ang } \\
& \left.\begin{array}{l}
\qquad 3 \\
\times \quad 2
\end{array}\right] \text { 分 } \\
& =2 \pi n \\
& -7 \leq x \leq-5 \\
& -s<x \\
& -(x+5)-(x+7)>1-(x+5)+x+7>1 \quad x+5+x+7>1 \\
& -x-5-x-7>1 \\
& -2 x-12>1 \\
& -2 x>13 \\
& x<-6.5 \\
& x<-7 \\
& \text { × } 6 \text { i no vies }
\end{aligned}
\]
\[
\begin{aligned}
& a, b \in(c, \alpha) \text { Gr } c<\alpha \text { Gr }|a-b|<\alpha-c \quad k \\
& \text { (כינו.) } \\
& a>b \text { o nij } \\
& a<\alpha \\
& -c>-b \quad \text { sk } c<b
\end{aligned}
\]
\[
\begin{array}{r}
+\left\{\begin{array}{r}
a<\alpha \\
-b<-c \\
a>b: a<\alpha-c \\
|a-b|<\alpha-c
\end{array}\right.
\end{array}
\]
\[
\begin{gathered}
0<a<b<1 \text { for } \quad|a+b| \leq 1 \\
a=0.5 \quad b=0.6 \quad \mid \supset) \quad \text { I.r } \\
(a+b)=1.1 \neq 1
\end{gathered}
\]
\[
\begin{gathered}
0<b<\alpha \quad \text { (sir } 0<a<c \quad \text { sr } \quad \frac{a}{b}<\frac{c}{d} \\
b=2 \quad b=1000 \quad a=1 \quad c=2 \quad \text { po) kr } \\
\frac{1}{2} \neq \frac{2}{1000}
\end{gathered}
\]
\[
\begin{array}{ll}
0<d<b \quad \text { or } 0<a<c \quad \frac{a}{b}<\frac{c}{d} \\
1>0 \\
\frac{a}{b}<\frac{c}{a} \\
a d<b c
\end{array}
\]
lre no nama ptk 5 \&, polla noon of dleens about st, lije abond if Viluma prem New AR
\(x_{1}, \ldots, x_{n} \in R\) res \(n \in N\) rs ：r．3．Ic
\[
\left|x_{1}+x_{2}-\ldots-x_{n}\right| \leq\left|x_{1}\right|+\left|x_{2}\right|+\ldots\left|x_{n}\right|
\]

נוכיח
i n
200 • 20 • \(01 \leq\left|x_{1}\right| \quad h=1\) －Fin lille the ins por \(\left|x_{1}-2 x_{2}\right| \leq\left|x_{1}\right|+\left|x_{2}\right| \quad n=2\)
－3 ÷ 12552
pissive hor socia in
\[
\left|x_{1}+\ldots+x_{k}\right| \leq\left|x_{1}\right|+\ldots+\left|x_{k}\right|
\]
行列 \(k+1\) 人）
\[
\begin{gathered}
\left|x_{1}+\ldots+x_{k}+x_{k+1}\right| \leq\left|x_{1}\right|+\ldots+\left|x_{k}\right|+\left|x_{k+1}\right| \\
x_{k \rightarrow 1}=B \quad x_{1}+\ldots+x_{k}=A \quad|-\infty| \\
|A+B| \leq|A|+|B|
\end{gathered}
\]
ank you is o jugy dina jinle kn
\[
\begin{aligned}
& |a-b| \geq||a|-|b||, \quad a, \lambda \in R \quad \text { rs }: r \cdot 3 \cdot a \\
& \text { |a| -> } \\
& |a|=|(a-b)+b| \quad: p<2 \quad \forall k \operatorname{sing}\left(r^{i}\right) \\
& \text { cren \| } \\
& |(a-b)+b| \leq|a-b|+|b| \\
& |a|-|b| \leq|a-b| \\
& \text { رحدرام د-|bi }
\end{aligned}
\]
\(|b|=|(b-a)+a|\) p R w/ ris
\[
|b| \leq|a-b|+|a|
\]
Slc
\[
{ }^{\nabla t}|a| \leadsto \text { iఎr }
\]
\[
\begin{aligned}
& |b|-|a| \leq|a-b| \\
& -|a-b| \leq|a|-|b| \\
& -|a-b| \leq|a|-|b| \leq|a-b| \\
& ||a|-|b|| \leq|a-b|
\end{aligned}
\]
\[
\begin{aligned}
& a, a_{0}, b, b_{0} \in R \quad \cdots{ }^{\prime} \\
& \left|a b-a_{0} b_{0}\right| \leq|a| \cdot\left|b-b_{0}\right| \pm\left|b_{0}\right| \cdot\left|a-a_{0}\right|:(3 \\
& \left|a b-a_{0} b \cdot\right|=\left|\left(a b-a b_{0}\right)+\left(a b_{0}-a_{0} b_{0}\right)\right|
\end{aligned}
\]
eliens prec \(k\) or
\[
\begin{gathered}
a b-a b_{0}-0 \times f\left(x_{0}\right) \\
a b_{0}-a_{0} b_{0}-0 y-1
\end{gathered}
\]
\[
\begin{aligned}
& \left|\left(a b-a b_{0}\right)+\left(a b_{0}-a_{0} b_{0}\right)\right| \leq\left|a b-a b_{0}\right|+\left|a b_{0}-a_{0} b_{0}\right| \\
& |=| \begin{array}{c}
\left|\left(a b-a b_{0}\right)+\left(a b_{0}-a_{0} b_{0}\right)\right| \leq\left|a \cdot\left(b-b_{0}\right)\right|+\left|b_{0} \cdot\left(a-a_{0}\right)\right| \\
|x \cdot y|=|x| \cdot|y| \\
\left|a b-a_{0} b_{0}\right| \leq|a \cdot| b-b_{0}\left|+\left|b_{0}\right| \cdot\right| a-a_{0} \mid
\end{array}
\end{aligned}
\]
دلاسَ (VIT) P.) Valulu a, b e c'jj anis. le prinkt)

bod fun shou neel as in inkly nan jee \(\left|a b-a_{0} b o\right|\) i.20 心. vien s'e
- lnk no win tor lea kar illa jo moil \(\left|b-b_{0}\right|,\left|a-a_{0}\right| \leq 0.01 s / c\) બૂગ 4 Iron Ima \(\left|a b-a b_{.}\right| \leq 10 \cdot 0.01+10 \cdot 0.01=0.2\) in 0.2 kn Non
\[
\begin{aligned}
& n \in N \text { os } \sum_{k=0}^{n} k^{2}=\frac{h \cdot(n+1) \cdot(2 n+1)}{6}: \Omega \cdot \varepsilon \\
& \text { ניכו כאשותצה } \\
& 0^{2}+1^{2}=1=\frac{1 \cdot(1+1) \cdot(2+1)}{6} \\
& 1 \stackrel{?}{=} \frac{1 \cdot 2 \cdot 3}{6} \\
& \text { - } 1=1
\end{aligned}
\]

Byrry 375
\[
\begin{aligned}
& \dot{\sum_{k=0}^{m} k^{2}=\frac{m \cdot(m+1) \cdot(2 m+1)}{6}} \\
& n=m+1 \text { infor argistar are in is } \\
& =\frac{(m+1) \cdot(m+2) \cdot(2 m+3)}{6} \text { is wirs }
\end{aligned}
\]

よこし
\[
\begin{aligned}
& \sum_{k=0}^{m+1} k^{2}=\sum_{k=0}^{n} k^{2}+(m+1)^{2}= \\
& \frac{m \cdot(m+1) \cdot(2 m+1)}{6}+(m+1)^{2}=\frac{r}{2} \underbrace{m+1}_{\text {are }} \text { (家) } \\
& \text { ( } \left.\frac{m \cdot(2 m+1)}{6}+(m+1)\right]=\left[\frac{2 m^{2}+m+6 \cdot(m+1)}{6}\right]=m+1 \cdot\left[\frac{2 m^{2}+7 m+6}{6}\right] \\
& \text { e) } \frac{(m+2) \cdot(2 m+3)}{6}
\end{aligned}
\]
\[
\begin{aligned}
& , a>-1 \quad \text { Gा } \quad n \in N \\
& (1+a)^{n} \geq 1+n a \quad \text {. } n=1 \\
& (1+a)^{1} \geq 1+1 \cdot a \\
& 1+a \geq 1+a
\end{aligned}
\]
✓ \(V\)

A3－12JŁD 283
\[
(1+a) \geq 1+k a
\]
\(n=k+1\) nor visus vaca is \(n>1\) is
\[
(1+a)^{k+1} \geq 1+(k+1) \cdot a
\]
： 5.3 N1／3
：Thup Toro
\[
\begin{aligned}
& (1+a)^{k+1}=(1+a)^{k} \cdot(1+a) \\
& (1+a)^{k} \geq 1+k a \quad \text { Ş}(12) \in(k)^{n .)} \text { or } \\
& (1+a)^{k} \cdot(1+a) \geq(1+k a) \cdot(1+a)
\end{aligned}
\]

マ／10 sinto sol N．
\[
(1+a)^{k} \cdot(1+a) \geq 1+(k+1) \cdot a \rightarrow k a^{2}
\]
we Hill wer espen of rie k \(a^{2}\)
\[
(1+a)^{k+1} \geq 1+(k+1) \cdot a
\]
\[
\begin{gathered}
a>0 \quad \text { rr } \quad n \in N \quad \text { r } \quad r \quad: \text { r.3 } \\
(1+a)^{n} \geq 1+n a+\frac{n \cdot(n-1)}{2} \cdot a^{2}
\end{gathered}
\]

\textbf{: 000 Min}
\[
\begin{array}{lll}
\checkmark & 1+a \geq 1+a+0 & h=1 \\
& (1+a)^{2} \\
& \stackrel{?}{\geq} 1+2 a+a^{2} & n=2 \\
\checkmark & (1-a)^{2}=1+2 a+a^{2} &
\end{array}
\]

\textbf{: \(3,12,1 k .28\)}
- iner \(k \geq 2\) nor no in
\[
(1+a)^{k} \geq 1+k a+\frac{k \cdot(k-1)}{2} \cdot a^{2}
\]

\[
\begin{aligned}
& \cdot(1+a)^{k+1} \geq 1+(k+1) \cdot a+\frac{(k+1) \cdot k}{2} \cdot a^{2} \\
& \left.(1+a)^{2} \geq 0, \cdots 3,12\right) k \cdot \min \geq 1 \\
& (1+a)^{k+1}=(1+a)^{k} \cdot(1+a) \geq\left(1+k a+\frac{k \cdot(k-1)}{2} a^{2}\right) \cdot(1+a) \\
& (1+a)^{k+1} \geq(1-1)+(k a \cdot 1)+\left(\frac{k \cdot(k-1)}{2} \cdot a^{2} \cdot 1\right)+ \\
& (1 \cdot a)+(k a \cdot a)+\left(\frac{k \cdot(k-1)}{2} \cdot a^{2} \cdot a\right) \\
& (1+a)^{k+1} \geq 1+k a+\frac{k \cdot(k-1)}{2} \cdot a^{2}+a+k a^{2}+ \\
& a^{3} \cdot \frac{k \cdot(k-1)}{2} \\
& (1+a)^{k+1} \geq 1+(k+1) \cdot a+\frac{k \cdot(k+1)}{2} \cdot a^{2}+\frac{k \cdot(k-1)}{2} a^{3}
\end{aligned}
\]
（ii）＇\(k\) fr uner 16 Nعت 3 tw pr \(k \geq 2\) ．or
\[
(1+a)^{k+1} \geqslant 1+(k+1) \cdot a+\frac{k \cdot(k+1)}{2} \cdot a^{2}
\]
erp
\[
\binom{n}{k}=\frac{n!}{k!(n-k)!} \quad x \in \operatorname{n} \in N \text { 万si } a, b \in k \quad G: r \cdot 3: 2
\]
\[
(a+b)^{n}=\sum_{k=0}^{n}\binom{n}{k} \cdot a^{k} \cdot b^{n-k^{p} ; k^{n}}
\]
－3，125－
\(n=1 \quad 0.00\) app
\[
\begin{aligned}
(a+b)^{1} \stackrel{?}{=} \sum_{k=0}^{1}\binom{1}{k} a^{k} b^{1-k}=\binom{1}{0} & -a^{0} \cdot b^{1} \\
& +\binom{1}{1} a^{1} b^{0}
\end{aligned}
\]
\[
\begin{gathered}
(a+b)^{1} \stackrel{?}{=}(1 \cdot 1 \cdot b)+(1 \cdot a \cdot 1) \\
a+b=b+a
\end{gathered}
\]
－
33，12生法 263
lice 5 of \(h=m\) ，
\[
(a+b)^{m}=\sum_{k=0}^{m}\binom{m}{k} a^{k} b^{m-k}
\]
\(m+1\) nor－
\[
(a+b)^{m+1}=(a+b)^{m} \cdot(a+b)=
\]
\[
\begin{aligned}
& {\left[\sum_{k=0}^{m}\binom{m}{k} a^{k} b^{m-k}\right] \cdot(a+b)={ }_{\left(r^{\prime} m\right)} r^{n}} \\
& a \cdot \sum_{k=0}^{m}\binom{m}{k} a^{k} b^{m-k}+b \cdot \sum_{k=0}^{m}\binom{m}{k} a^{k} b^{m-k}= \\
& \sum_{k=0}^{m}\binom{m}{k} \cdot a^{k+1} \cdot b^{m-k}+\sum_{k=0}^{m}\binom{m}{k} \cdot a^{k} \cdot b^{m+1-k} \\
& \text { REIDO re yesi } j=k+1 \text { yes } \\
& \sum_{j}^{m+1}\binom{m}{j-1} a^{j} b^{m+1-j}+\sum_{k=0}^{m}\binom{m}{k} \cdot a^{k} \cdot b^{m+1-k}= \\
& \sum_{k}^{m+1}\binom{m}{k-1} a^{k} b^{m+1-k}+\sum_{k=0}^{m}\binom{m}{k} a^{k} b^{m+1-k}= \\
& a^{m+1}+b^{m+1} \rightarrow \sum_{k=1}^{m}\left[\binom{m}{k-1}+\binom{m}{k}\right] a^{k} b^{m+1-k} \\
& \binom{m}{k-1}+\binom{m}{k}=\binom{m+1}{k} \\
& (a+b)^{m+1}=a^{m+1}+b^{m+1}+\sum_{k=1}^{n}\binom{m+1}{k} a^{k} b^{m+1-k} \\
& (m+1): L^{m+1-0}=b^{m+1} \quad k=0 \quad x \in D
\end{aligned}
\]
\[
\binom{m+1}{m+1} a^{m+1} b^{m+1-(m+1)}=a^{m+1} \quad k=m+1 \quad x<0
\]
\(\left.\begin{array}{r}\text { Incol } \\ \text { - " }(x) " \\ \text { spl }\end{array} \right\rvert\, x\)
\[
=\sum_{k=0}^{m+1}\binom{m+1}{k} a^{k} b^{m+1-k}
\]
\textbf{(Q)}
N
\[
\begin{aligned}
& \prod_{k=2}^{n}\left(1-\frac{1}{k^{2}}\right)=\frac{n+1}{2 n} \\
& \quad \because 3,12, \cdots \infty
\end{aligned}
\]
\(n=2\) oo life
\[
\begin{array}{rlr}
\prod_{k=2}^{2}\left(1-\frac{1}{k^{2}}\right)=1-\frac{1}{2^{2}}=1-\frac{1}{4} & =\frac{3}{4} & \text { :Due Tek } \\
& \frac{2+1}{2 \cdot 2}=\frac{3}{4} \quad \text { ivi Tck }
\end{array}
\]
- evas lille rpro

ミ3,12J.k. 263
\(2 \leq m\). Ixfs حد \(m\), Nore \(m\), Jose
\[
\prod^{m}\left(1-\frac{1}{k^{2}}\right)=\frac{m+1}{2 m}
\]

2-2
\(m+1\) inivi
\[
\begin{aligned}
\prod_{k=2}^{m+1} & \left(1-\frac{1}{k^{2}}\right)=\left[\prod_{k=2}^{m}\left(1-\frac{1}{k^{2}}\right)\right] \cdot\left(1-\frac{1}{(m+1} .\right. \\
& \frac{m+1}{2 m} \cdot\left(1-\frac{1}{(m+1)^{2}}\right)= \\
& \frac{m+1}{2 m} \cdot\left(\frac{(m+1)^{2}-1}{(m+1)^{2}}\right)= \\
& \frac{m+1}{2 m} \cdot \frac{m^{2}+2 m}{(m+1)^{2}}= \\
& \frac{m+1}{2 m} \cdot \frac{m \cdot(m+2)}{(m+1)^{2}}=\frac{m+2}{2 \cdot(m+1)}= \\
& \frac{(m+1)+1}{2 \cdot(m+1)}
\end{aligned}
\]

\[
\begin{aligned}
& 2 \leq n \in N \text { r.r :r.3 } \quad \begin{array}{r}
1 \\
\frac{1}{n+1}+\frac{1}{n+2}+\ldots+\frac{1}{2 n}>\frac{13}{24}
\end{array}
\end{aligned}
\]

Pinivies -es non \(n=2 \quad 0.00\) injer
\[
\frac{1}{2}+\frac{1}{2}=\frac{4+3}{2}=7 \cdot \frac{14}{1} \text { : Rove Terc }
\]
\[
\frac{14}{24}>\frac{13}{24}
\]
:
\(n=k\) is
\[
S_{k}=\frac{1}{k+1}+\frac{1}{k+2}+\ldots+\frac{1}{2 k}>\frac{13}{24}
\]
: [. 3 ws. \(h=k+1\) wor we now iss]
\[
S_{k \rightarrow 1}=\frac{1}{k \rightarrow 2}+\ldots+\frac{1}{2 k}+\frac{1}{2 k \rightarrow 1}+\frac{1}{2 k+2}>\frac{13}{24}
\]
\[
S_{k+1}=S_{k}-\frac{1}{k+1}+\frac{1}{2 k+1}+\frac{1}{2 k+2}
\]
of lexk ibou wer fille kile night ro
\[
\begin{aligned}
& \frac{1}{2 k+1}+\frac{1}{2 k+2}-\frac{1}{k+1} \geq 0 \\
& \frac{1}{k+1}+\frac{1}{2 \cdot(k+1)}-\frac{1}{k+1}= \\
& \text { 1) }-\frac{1}{k+1}=\frac{1}{2 \cdot(k+1)}-\frac{2}{2 \cdot(k+1)}=-\frac{1}{2 \cdot(k+1)}
\end{aligned}
\]
\(r_{\text {appl }}\) יוֹא
\[
\frac{1}{2 k+1}-\frac{1}{2 \cdot(k+1)}
\]

Gpl Eien von ron
\[
\begin{aligned}
& \frac{1 \cdot[2 \cdot(k+1)]-1 \cdot[2 k+1]}{(2 k+1) \cdot 2 \cdot(k+1)}= \\
& \frac{2 k+2-2 k-1}{(2 k+1) \cdot(2 k+2)}=\frac{1}{(2 k+1)(2 k+2)}>0 \\
& (2-1 \text { or } 2 \text { in } k+0)
\end{aligned}
\]

\end{document}
